\subsection{معرفی و انگیزه}
این ماژول در چارچوب پژوهش حاضر پیاده‌سازی و ارزیابی شده و فاقد مقاله مستقل جداگانه در ادبیات تحقیق است.

در تقریب رتبه پایین وزن‌های ترنسفورمر به وسیله تجزیه مقادیر منفرد (\lr{SVD})، انتخاب رتبه بریده‌شده $r$ یکی از چالش‌های بنیادین است. در رویکردهای سنتی، یک رتبه ثابت و یکسان برای تمامی لایه‌ها اعمال می‌شود؛ این در حالی است که ماهیت ریاضی و پراکندگی مقادیر منفرد در لایه‌های مختلف به شدت ناهمگن است. برخی لایه‌ها دارای طیف مقادیر منفرد بسیار چوله‌دار با زوال سریع هستند (که با رتبه‌های بسیار پایین با دقت بالا قابل تقریب‌اند)، در حالی که لایه‌های دیگر دارای توزیع یکنواخت انرژی بوده و نیازمند حفظ ابعاد بالاتری هستند.

ماژول \lr{MGIE-SVD (Maximum Information Entropy SVD)} بر پایه اصل بیشینه‌سازی آنتروپی اطلاعاتی مقادیر منفرد طراحی شده است. این روش با محاسبه آنتروپی شانون برای توزیع احتمالاتی مقادیر منفرد در هر لایه، میزان پراکندگی طیفی لایه را به صورت خودکار ارزیابی نموده و رتبه تقریب را متناسب با آنتروپی اطلاعاتی آن لایه تخصیص می‌دهد. این امر موجب می‌شود لایه‌هایی که محتوای اطلاعاتی پیچیده‌تر و طیف یکنواخت‌تری دارند رتبه بالاتری دریافت کنند و لایه‌های با مقادیر منفرد غالب با رتبه‌های فشرده‌تر تقریب زده شوند.
